% Section 16.5: Exact and closed $1$-forms: the first Stokes theorem | Part: 3 | Depends on: sec:fields, sec:tangent | Status: done | Original lines: 2538-2576
%!TEX root = ../../main.tex
\subsection{Exact and closed $1$-forms: the first Stokes theorem}\label{sec:expl-closed}

For a function $f$ the differential is the $1$-form $\dd f=\partial_if\,\dd x^i$, so $\langle\dd f,v\rangle=v^i\partial_if=v(f)$ (the directional derivative). Along a curve:
\begin{align*}
\int_\gamma\dd f&=\int_a^b\partial_if(\gamma(t))\,\dot\gamma^i(t)\,\dd t &&\text{(definition \eqref{eq:lineint})}\\
&=\int_a^b\frac{\dd}{\dd t}f(\gamma(t))\,\dd t &&\text{(chain rule)}\\
&=f(\gamma(b))-f(\gamma(a)) &&\text{(fundamental theorem of calculus).}
\end{align*}
\begin{equation}\label{eq:ftc}
\int_\gamma\dd f=f(\gamma(b))-f(\gamma(a)).
\end{equation}
This is Stokes' theorem for $k=1$: the integral of $\dd f$ over the curve equals $f$ evaluated on the boundary $\partial\gamma=\{\gamma(b)\}-\{\gamma(a)\}$.

A form is \emph{exact} if $\omega=\dd f$ and \emph{closed} if $\dd\omega=0$.
\begin{itemize}
\item \emph{Exact $\Rightarrow$ closed.} $\dd(\dd f)=\dd(\partial_if\,\dd x^i)=\partial_j\partial_if\,\dd x^j\wedge\dd x^i=0$, because $\partial_j\partial_if$ is symmetric in $(i,j)$ while $\dd x^j\wedge\dd x^i$ is antisymmetric.
\item \emph{Exact $\Rightarrow$ integral zero on closed loops.} If $\gamma(b)=\gamma(a)$, \eqref{eq:ftc} gives $0$.
\item The converse ``closed $\Rightarrow$ exact'' fails in general: this is the starting point of de Rham cohomology (next example).
\end{itemize}

\subhead{How to compute $\dd$ of a $1$-form in the plane}
For $\omega=P\,\dd x+Q\,\dd y$:
\begin{align*}
\dd\omega&=\dd P\wedge\dd x+\dd Q\wedge\dd y\\
&=(P_x\dd x+P_y\dd y)\wedge\dd x+(Q_x\dd x+Q_y\dd y)\wedge\dd y\\
&=P_y\,\dd y\wedge\dd x+Q_x\,\dd x\wedge\dd y=(Q_x-P_y)\,\dd x\wedge\dd y,
\end{align*}
using $\dd x\wedge\dd x=\dd y\wedge\dd y=0$ and $\dd y\wedge\dd x=-\dd x\wedge\dd y$.

\begin{example}[The angle form]\label{ex:angle}
On $M=\R^2\setminus\{0\}$ let $\omega=\dfrac{-y\,\dd x+x\,\dd y}{x^2+y^2}=P\,\dd x+Q\,\dd y$, $P=\dfrac{-y}{x^2+y^2}$, $Q=\dfrac x{x^2+y^2}$.
\begin{description}[leftmargin=0pt,labelindent=0pt,font=\bfseries]
\item[Step 1: $\omega$ is closed.] With $r^2=x^2+y^2$: $Q_x=\dfrac{r^2-x\cdot2x}{r^4}=\dfrac{y^2-x^2}{r^4}$, $P_y=-\dfrac{r^2-y\cdot2y}{r^4}=\dfrac{y^2-x^2}{r^4}$, so $\dd\omega=(Q_x-P_y)\,\dd x\wedge\dd y=0$.
\item[Step 2: why it is called ``the angle''.] With $\theta=\arctan(y/x)$, $\partial_x\theta=\dfrac{-y/x^2}{1+y^2/x^2}=\dfrac{-y}{x^2+y^2}$ and $\partial_y\theta=\dfrac{1/x}{1+y^2/x^2}=\dfrac x{x^2+y^2}$, so $\dd\theta=\omega$ locally. (The function $\theta$ is multivalued around the origin.)
\item[Step 3: integrate on the unit circle] $\gamma(t)=(\cos t,\sin t)$, $0\le t\le2\pi$: velocity $\dot\gamma=(-\sin t,\cos t)$; on the circle $r=1$, $\omega_x=P=-\sin t$, $\omega_y=Q=\cos t$; contraction $\langle\omega_{\gamma(t)},\dot\gamma(t)\rangle=(-\sin t)(-\sin t)+(\cos t)(\cos t)=1$; integral $\oint_\gamma\omega=\int_0^{2\pi}1\,\dd t=2\pi$.
\item[Step 4: conclusion.] If $\omega=\dd f$ with $f$ single-valued on $M$, then \eqref{eq:ftc} would give $\oint_\gamma\omega=f(\gamma(2\pi))-f(\gamma(0))=0$, contradicting Step 3. Hence $\omega$ is closed but not exact on $M$: there is no global angle function. The number $\frac1{2\pi}\oint\omega$ is the \emph{winding number}, and $H^1_{\mathrm{dR}}(\R^2\setminus\{0\})\cong\R$ \cite{BottTu}.
\end{description}
\end{example}
